\documentclass[12pt, openright, a4paper, twoside, english, brazil]{article}
\usepackage[top=2cm, bottom=2cm, left=2cm, right=2cm]{geometry}

\usepackage{mathtools}
\usepackage{amsfonts}
\usepackage{amssymb}
\usepackage{amsmath}
\usepackage[utf8]{inputenc}
\usepackage[portuguese]{babel}
\usepackage[T1]{fontenc}
\usepackage{titlesec}

\usepackage{libertine}
\usepackage{newtxmath}

\usepackage{enumerate}
\usepackage{tikz}
\usepackage{pgfplots}
\usepackage{hyperref}
\hypersetup{
  colorlinks=false,
  urlbordercolor=red,
  pdfborderstyle={/S/U/W 1}
}

\newenvironment{solution}{\color{red}\noindent\textbf{Solução:}\\}{\color{black}}
\newcommand{\df}[2]{\frac{d#1}{d#2}}
\newcommand{\ghlink}[2]{\href{https://github.com/dreisss/study-notes/blob/main/#1}{#2}}
\newcommand{\difflink}[1]{\ghlink{Differentiation/Real/Abstract_pt.pdf}{#1}}

\titleformat{\section}[block]{\bf\raggedright\normalsize}{\thesection}{1em}{}[]
\titlespacing*{\section}{0pt}{1.5em}{0.75em}[0pt]
\titleformat{\subsection}[block]{\raggedright\normalsize}{\bf\thesubsection.}{1em}{}[]
\titlespacing*{\subsection}{0pt}{1.5em}{0.75em}[0pt]

% -----------------------------------------------------------------------------

\title{Soluções - Redes Neurais Artificiais: Princípios e prática}
\author{Capítulo 1 - Introdução}
\date{}

\begin{document}

\maketitle

\section{Modelos de um neurônio}

{\bf Pré-requisitos:}
\ghlink{Functions/Abstract_pt.pdf}{Funções},
\ghlink{Functions/Abstract_pt.pdf}{Limites*},
\difflink{Derivadas}

\subsection{Um exemplo de função logística é definida por}
$$
\varphi(v)=\frac{1}{1+\exp(-av)}
$$
cujos valores limites são $0$ e $1$. Mostre que a derivada de
$\varphi(v)$ em relação a $v$ é dada por
$$
\df{\varphi}{v}=a\varphi(v)[1-\varphi(v)]
$$
Qual é o valor desta derivada na origem?\\

\begin{solution}

  Para resolver podemos utilizar a \difflink{regra da cadeia}. Tomando
  a substituição com $u(v)=1+e^{-av}$, sabemos que $u'(x)=-ae^{-av}$.
  Portanto temos:
  $$\varphi(v)=\frac{1}{u(v)}=u^{-1}(v)\implies\varphi'(v)=-u^{-2}(v)u'(v)$$
  $$\varphi'(v)=\frac{ae^{-av}}{1+e^{-av}}=ae^{-av}\varphi^2(v)$$
  Podemos isolar $e^{-av}$ em função de $\varphi(v)$:
  $$e^{-av}=\frac{1}{\varphi(v)}-1=\frac{1-\varphi(v)}{\varphi(v)}$$
  E por fim:
  $$
  \df{\varphi}{v}(v)=\frac{a\varphi^2(v)[1-\varphi(v)]}{\varphi(v)}\implies
  \df{\varphi}{v}(v)=a\varphi(v)[1-\varphi(v)]
  $$

  Para calcular o valor da derivada na origem basta calcular:
  $\displaystyle \varphi(0)=\frac{1}{1+e^{-a\cdot0}}=\frac{1}{2}$. E
  aplicando na derivada:
  $$
  \df{\varphi}{v}(0)=a\varphi(0)[1-\varphi(0)]\implies
  \df{\varphi}{v}(0)=\frac{a}{4}
  $$

  Para criar um esboço da função podemos tomar os seguintes limites:
  $$
  \lim_{v\to+\infty}\varphi(v)=\lim_{v\to+\infty}\frac{1}{1+e^{-av}},
  \lim_{v\to+\infty}e^{-\infty}=0\implies \lim_{v\to+\infty}\varphi(v)=1
  $$
  $$
  \lim_{v\to-\infty}\varphi(v)=\lim_{v\to-\infty}\frac{1}{1+e^{-av}},
  \lim_{v\to-\infty}e^{+\infty}=\infty\implies \lim_{v\to+\infty}\varphi(v)=0
  $$

  \newpage

  Portanto, o gráfico da função é:
  \begin{center}
    \begin{tikzpicture}
      \begin{axis}[
          axis lines = center,
          xlabel = \(v\),ylabel ={\(\varphi(v)\)},
          yscale = 0.5,
          xtick distance = 11, ytick distance = 1,
          extra x ticks={0}, extra x tick labels={0},
          extra y ticks={0.5}, extra y tick labels={\(\frac{a}{2}\)},
        ]
        \addplot [domain=-10:10,color=red]{1 / (1 + exp(-x))};
      \end{axis}
    \end{tikzpicture}
  \end{center}

\end{solution}

\subsection{Uma função sigmóide ímpar é definida por}
$$
\varphi(v)=\frac{1-\exp(-av)}{1+\exp(-av)}=\tanh(\frac{av}{2})
$$
onde $\tanh$ representa a tangente hiperbólica. Os valores limites
desta segunda função sigmóide são $-1$ e $+1$. Mostre que a derivada de
$\varphi(v)$ em relação a $v$ é dada por
$$
\df{\varphi}{v}=\frac a2[1-\varphi^2(v)]
$$
Qual é  valor desta derivada na origem? Suponha que o parâmetro de
inclinação $a$ seja infinitamente grande. Qual é a forma resultante de
$\varphi(v)$?\\

\begin{solution}

  Podemos resolver utilizando a \difflink{derivada da divisão}, a
  \difflink{regra da cadeia} e o
  \ghlink{/BasicAlgebra/SpecialProducts.pdf}{quadrado da soma}.
  Aplicando temos:
  $$\varphi'(v)=\frac{ae^{-av}(1+e^{-av})+ae^{-av}(1-e^{-av})}{(1+e^{-av})^2}=a\frac{a[]}{(1+e^{-av})^2}$$

\end{solution}

\subsection{Uma outra função sigmóide ímpar é a sigmóide algébrica:}
$$
\varphi(v)=\frac v{\sqrt{1+v^2}}
$$
cujos valores limites são $-1$ e $+1$. Mostre que a derivada de
$\varphi(v)$ em relação a $v$ é dada por
$$
\df{\varphi}{v}=\frac{\varphi^3(v)}{v^3}
$$
Qual é o valor desta derivada na origem?

\subsection{Considere as duas seguintes funções:}
\begin{enumerate}[\bfseries (i)]
  \item
    $\displaystyle\varphi(v)=\frac1{\sqrt{2\pi}}\int_{-\infty}^{\infty}\exp(-\frac{x^2}{2})dx$

  \item $\displaystyle\varphi(v)=\frac2\pi\tan^{-1}(v)$
\end{enumerate}
Explique por que estas funções satisfazem os requisitos de uma
função sigmóide. De que modo estas duas funções diferem entre si?

\subsection{Qual das cinco funções sigmóides descritas nos Problemas
  1.1 a 1.4 seria qualificada como uma função de distribuição (de
  probabilidade) cumulativa? Justifique a sua resposta.
}

\subsection{Considere a função de ativação pseudolinear $\varphi(v)$
mostrada na Fig. P1.6.}
\begin{center}
  \begin{tikzpicture}
    \begin{axis}[
        axis lines = center,
        xlabel = \(v\),ylabel = \(\varphi(v)\),
        yscale = 0.5,
        xtick distance=3, ytick distance=3,
        xmin=-2, xmax=2,
        ymin=0, ymax=2,
        legend pos=north west,
        extra x ticks={0, -0.5, 0.5},
        extra x tick labels={0, \(-0.5a\), \(0.5a\)},
        extra y ticks={0.5, 1},
        extra y tick labels={\(0.5b\), \(b\)}
      ]

      \addplot [domain=-5:5, color=blue] coordinates {(-5, 0)(-0.5,0)};
      \addplot [domain=-0.5:0.5, color=blue] coordinates {(-0.5, 0)(0.5,1)};
      \addplot [domain=-5:5, color=blue] coordinates {(0.5,1)(5,1)};

      \addplot [domain=-5:5, style=dashed] coordinates {(0.5,0)(0.5,1)};
      \addplot [domain=-5:5, style=dashed] coordinates {(0,1)(0.5,1)};
    \end{axis}
  \end{tikzpicture}
\end{center}
\begin{enumerate}[a)]
\item Formule $\varphi(v)$ como uma função de $v$.
\item O que acontece com $\varphi(v)$ se $a$ aproximar-se de zero?
\end{enumerate}

\subsection{Repita o Problema 1.6 para a função de ativação
pseudolinear $\varphi(v)$ mostrada na Fig. P1.7.}
\begin{center}
\begin{tikzpicture}
  \begin{axis}[
      axis lines = center,
      xlabel = \(v\),ylabel = \(\varphi(v)\),
      xtick distance=3, ytick distance=3,
      xmin=-2, xmax=2,
      ymin=-2, ymax=2,
      legend pos=north west,
      extra x ticks={-1, 0, 1},
      extra x tick labels={\(-a\), 0, \(a\)},
      extra y ticks={-1, 1},
      extra y tick labels={\(-b\), \(b\)}
    ]

    \addplot [domain=-5:5, color=blue] coordinates {(-5, -1)(-1,-1)};
    \addplot [domain=-0.5:0.5, color=blue] coordinates {(-1, -1)(1,1)};
    \addplot [domain=-5:5, color=blue] coordinates {(1,1)(5,1)};

    \addplot [domain=-5:5, style=dashed] coordinates {(1,0)(1,1)};
    \addplot [domain=-5:5, style=dashed] coordinates {(0,1)(1,1)};
    \addplot [domain=-5:5, style=dashed] coordinates {(-1,0)(-1,-1)};
    \addplot [domain=-5:5, style=dashed] coordinates {(0,-1)(-1,-1)};
  \end{axis}
\end{tikzpicture}
\end{center}

\subsection{Um neurônio tem uma função de ativação $\varphi(v)$
definida pela função logística do problema 1.1, onde $v$ é o campo
local induzido, e o parâmetro de inclinação $a$ está disponível para
ajustes. Considere que $x_1,x_2,...,x_m$, representem os sinais de
entrada aplicados aos nós de fonte do neurônio e que $b$ represente o
bias. Por conveniência de representação, podemos fazer com que o
parâmetro de inclinação $a$ seja absorvido pelo campo local induzido
$v$, escrevendo}
$$
\varphi(v)=\frac1{1+\exp(-v)}
$$
Como você modificaria as entradas $x_1,x_2,...,x_m$ de forma a
produzir a mesma saída de antes? Justifique a sua resposta.

\subsection{Um neurônio $j$ recebe entradas de quatro outros
neurônios cujos níveis de ativação são $10, -20,4,-2$. Os respectivos
pesos sinápticos do neurônio $j$ são $0.8, 0.2, -1.0$ e $-0.9$.
Calcule a saída do neurônio $j$ para as duas seguintes situações:}
\begin{enumerate}[(a)]
\item O neurônio é linear
\item O neurônio é representado por um modelo de McCulloch-Pitts.
\end{enumerate}
Assuma que o bias aplicado ao neurônio é zero.

\subsection{Repita o Problema 1.9 para um modelo de neurônio baseado
na função logística}
$$
\varphi(v)=\frac 1{1+\exp(-v)}
$$

\subsection{}
\begin{enumerate}[(a)]
\item Mostre que o modelo formal de McCulloch-Pitts de um neurônio
  pode ser aproximado por um neurônio sigmóide (i.e., um neurônio que
  utiliza uma função de ativação sigmóide) com pesos sinápticos grandes.
\item Mostre que o neurônio linear pode ser aproximado por um
  neurônio sigmóide com pesos sinápticos pequenos.
\end{enumerate}

\section{Arquitetura de rede}
\setcounter{section}{1}
\setcounter{subsection}{11}

\subsection{Uma rede alimentada adiante totalmente conectada tem 10
nós de fonte, 2 camadas ocultas, uma com 4 neurônios e a outra com 3
neurônios em um único neurônio de saída. Construa um grafo
arquitetural desta rede.}

\subsection{}
\begin{center}
\begin{tikzpicture}[shorten >=1pt,->]
  \tikzstyle{vertex}=[thick, circle, draw, inner sep = 1.5pt]

  \node[vertex] (N_1_1) at (0, 1) {};
  \node[vertex] (N_1_2) at (0,-1) {};
  \node[vertex] (N_2_1) at (2, 1) {};
  \node[vertex] (N_2_2) at (2,-1) {};
  \node[vertex] (N_3_1) at (4, 1) {};
  \node[vertex] (N_3_2) at (4,-1) {};
  \node[vertex] (N_4_1) at (6, 1) {};
  \node[vertex] (N_4_2) at (6,-1) {};
  \node[vertex] (N_5_1) at (8, 1) {};
  \node[vertex] (N_5_2) at (8,-1) {};
  \node[vertex] (N_6_1) at (10,0) {};
  \node[vertex] (N_7_1) at (12,0) {};

  \draw [thick, ->] (N_1_1) -- (N_2_1) node[midway, above] {5};
  \draw [thick, ->] (N_1_2) -- (N_2_2) node[midway, below] {-3};
  \draw [thick, ->] (N_1_1) -- (N_2_2) node[pos=0.65, below ] {2};
  \draw [thick, ->] (N_1_2) -- (N_2_1) node[pos=0.65, above ] {1};

  \draw [thick, ->] (N_2_1) -- (N_3_1) node[midway, above] {\(\varphi(\cdot)\)};
  \draw [thick, ->] (N_2_2) -- (N_3_2) node[midway, below] {\(\varphi(\cdot)\)};

  \draw [thick, ->] (N_3_1) -- (N_4_1) node[midway, above] {3};
  \draw [thick, ->] (N_3_2) -- (N_4_2) node[midway, below] {6};
  \draw [thick, ->] (N_3_1) -- (N_4_2) node[pos=0.65, below] {4};
  \draw [thick, ->] (N_3_2) -- (N_4_1) node[pos=0.65, above] {-1};

  \draw [thick, ->] (N_4_1) -- (N_5_1) node[midway, above] {\(\varphi(\cdot)\)};
  \draw [thick, ->] (N_4_2) -- (N_5_2) node[midway, below] {\(\varphi(\cdot)\)};

  \draw [thick, ->] (N_5_1) -- (N_6_1) node[midway, above] {-2};
  \draw [thick, ->] (N_5_2) -- (N_6_1) node[midway, below] {1};

  \draw [thick, ->] (N_6_1) -- (N_7_1) node[midway, above] {\(\varphi(\cdot)\)};
\end{tikzpicture}
\end{center}

\begin{enumerate}[(a)]
\item A Figura 1.13 mostra um grafo de fluxo de sinal de uma rede
  2-2-2-1 alimentada adiante. A função $\varphi(\cdot)$ representa
  uma função logística. Escreva o mapeamento de entrada-saída
  definido por esta rede.
\item Suponha que o neurônio de saída do grafo de fluxo de sinal da
  Fig. P1.13 opere na sua região linear. Escreva o mapeamento de
  entrada-saída definido por esta nova rede.
\end{enumerate}

\subsection{A rede descrita na Fig. P1.13 não tem bias. Suponha que
bias iguais a $-1$ e $+1$ sejam aplicados aos neurônios superior e
inferior da primeira camada oculta, e bias iguais a $+1$ e $-2$ sejam
aplicados aos neurônios superior e inferior da segunda camada oculta,
respectivamente. Escreva a nova forma do mapeamento de entrada-saída
definido pela rede.}

\subsection{Considere uma rede de múltiplas camadas alimentada
adiante, na qual todos os neurônios operam nas suas regiões lineares.
Justifique a afirmação de que esta rede é equivalente a uma rede
alimentada adiante de camada única.}

\subsection{Construa uma rede totalmente recorrente com 5 neurônios,
mas sem auto-realimentação.}

\subsection{A Figura P1.17 mostra um grafo de fluxo de sinal de uma
rede recorrente constituída de dois neurônios. Escreva a equação de
diferenças não-linear que define a evolução de $x_1(n)$ ou de
$x_2(n)$. Estas duas variáveis definem as saídas dos neurônios
superior e inferior, respectivamente. Qual é a ordem desta equação?}
\begin{center}
\begin{tikzpicture}[shorten >=1pt,->]
  \tikzstyle{neuron}=[circle, fill=black!20, draw, inner sep = 7pt]
  \tikzstyle{delay}=[rectangle, draw, minimum size = 0.8cm, minimum
  height = 0.8cm]

  \node[delay] (D_1) at (0, 3) {\(z^{-1}\)};
  \node[delay] (D_2) at (0, 2) {\(z^{-1}\)};
  \node[neuron] (N_1) at (0, 1) {};
  \node[neuron] (N_2) at (0, 0) {};

  % alternativa com curvas suaves
  % \draw [thick, ->]  (N_1) to [out=0, in=0 ] (D_1) node[] {};
  % \draw [thick, ->] (N_2) to [out=0, in=0] (D_2) node[] {};
  % \draw [thick, ->] (D_1) to [out=180, in=180] (N_2) node[] {};
  % \draw [thick, ->] (D_2) to [out=180, in=180] (N_1) node[] {};
  %
  \draw [thick, ->] (N_1) -- ++(1,0) |- (D_1);
  \draw [thick, ->] (N_2) -- ++(1.5,0) |- (D_2);

  \draw [thick, ->] (D_1) -- ++(-1.5,0) |- (N_2);
  \draw [thick, ->] (D_2) -- ++(-1,0) |- (N_1);
\end{tikzpicture}
\end{center}

\subsection{A Figura P1.18 mostra o grafo de fluxo de sinal de uma
rede recorrente que consiste de dois neurônios com
auto-realimentação. Escreva o sistema acoplado de duas equações de
diferenças de primeira ordem que descrevem a operação do sistema.}
\begin{center}
\begin{tikzpicture}[shorten >=1pt,->]
  \tikzstyle{neuron}=[circle, fill=black!20, draw, inner sep = 7pt]
  \tikzstyle{delay}=[rectangle, draw, minimum size = 0.8cm, minimum
  height = 0.8cm]

  \node[delay] (D_1) at (0, 3) {\(z^{-1}\)};
  \node[delay] (D_2) at (0, 2) {\(z^{-1}\)};
  \node[neuron] (N_1) at (0, 1) {};
  \node[neuron] (N_2) at (0, 0) {};

  \draw [thick, ->] (N_1) -- ++(1,0) |- (D_1);
  \draw [thick, ->] (N_2) -- ++(1.5,0) |- (D_2);

  \draw [thick, ->] (D_1) -- ++(-1.5,0) -- ++(0, -3.5) -- (N_2);
  \draw [thick, ->] (D_1) -- ++(-1.5,0) -- ++(0,-1.5) -- (N_1);
  \draw [thick, ->] (D_2) -- ++(-1,0) -- ++(0,-1.25) -- (N_1);
  \draw [thick, ->] (D_2) -- ++(-1,0) -- ++(0,-1.75) -- (N_2);

\end{tikzpicture}
\end{center}

\subsection{Uma rede recorrente tem 3 nós de fonte, 2 neurônios
ocultos e 4 neurônios de saída. Construa um grafo arquitetural que
descreva esta rede.}

\section{Representação do conhecimento}
\setcounter{section}{1}
\setcounter{subsection}{19}

\subsection{Uma forma útil de pré-processamento se baseia no modelo
auto-regressivo (AR) descrito pela equação de diferenças (para dados
de valores reais)}
$$
y(n)=w_1y(x-1)+w_2y(n-2)+...+w_My(n-M)+v(n)
$$
onde $y(n)$ é a saída do modelo; $v(n)$ é uma amostra retirada de um
processo de ruído branco com média zero e uma variância predefinida;
$w_1,w_2,...,w_M$, são os coeficientes do modelo AR; e M é a ordem do
modelo. Mostre que o uso deste modelo fornece duas formas de
invariância geométrica: (a) em escala e (b) em translação temporal.
Como estas duas invariâncias poderiam ser utilizadas em redes neurais?

\subsection{Considere que $\mathbf x$ seja um vetor de entrada e que
$\mathbf s(\alpha,\mathbf x)$ seja um operador de transformação agindo
sobre $\mathbf x$ e dependente de um parâmetro $\alpha$. O operador
$\mathbf s(\alpha, \mathbf x)$ satisfaz dois requisitos:}
\begin{itemize}
\item $\mathbf s(0, \mathbf x)=0$
\item $\mathbf s(0, \mathbf x)$ é diferenciável em relação a $\alpha$.
\end{itemize}
O vetor tangente é definido pela derivada parcial $\partial\mathbf
s(\alpha,\mathbf x)/\partial\alpha$ (Simard et al., 1992). Suponha
que $\mathbf x$ represente uma imagem e que $\alpha$ seja um
parâmetro de rotação. Como você calcularia o vetor tangente para o
caso em que $\alpha$ é pequeno? O vetor tangente é localmente
invariante em relação à rotação da imagem original; por quê?

\end{document}
